\section{Discourse Analysis}
\label{sec:discourse-analysis}

As a core conceptual part of this thesis, discourse analysis provides a theoretical foundation for examining how language is used to construct different meanings in competitive environments. This is not only limited to the area of Artificial Intelligence, but also to the field of linguistics and communication studies, although this research will be focused on the use of Natural Language Processing techniques.

\subsection{Discourse Analysis in Competitive Contexts}

A key objective of discourse analysis, and present research, is to identify the key linguistic features of each discourse. As mentioned before, there is particular interest in identifying these features in the context of competitive scenarios. Since the area of esports broadcasting is a relatively recent domain, there is a clear gap in the literature when it comes to discourse analysis in this context. However, traditional sports do have some investigation in this area, and by reviewing this literature, we can identify some key features that may be relevant to our research.

\medskip

Large-scale sporting events have been, therefore, crucial for discourse analysis, as it is expected to have experienced commentators and high-quality broadcasts. This was the case for the FIFA World Cup 2022, analyzed by Syaputri et al. \cite{syaputri2024charla}, with a focus on the different techniques used by commentators to talk about the participating teams. 

\medskip

Their findings are relevant for our research, as they provide a starting point for identifying some of the key properties of esports commentary, such as:

\begin{itemize}
  \item Changes in \textbf{register} between formal and informal talk depending on the state of the game. 
  \item Construction of identity through the use of the first person plural, to create an engaging narrative.
  \item Use of politeness strategies to maintain credibility and objectiveness when performing analysis.
\end{itemize}

However, despite its relevance, it must be noted that this article was retracted following concerns regarding publication manipulation by its journal. Therefore, their findings must be taken only as a guide - due to the overlap of properties and labels -, and not as a solid basis for our research. Furthermore, this reflects the lack of robust literature when it comes to discourse research, and the need for a proper annotated dataset which is a key contribution of this thesis. 

\medskip

More specific discourse properties, such as indicators of intensity in commentary, have been analyzed in order to identify key moments in games through machine learning techniques. These approaches overlap in objectives with the present research; enhancing the viewer experience, and lowering the cognitive load of broadcasters, and external spectators. 

\medskip 

A key article which also enters the Artificial Intelligence field is the work of Oliveira et al. \cite{oliveira2026multi} which explores the application of these techniques to both the context of traditional sports and esports. In this case, the use of multi-modal data, such as audio and video, is used to identify key moments in games, to highlight them for spectators. In a similar fashion to the present thesis, part of the goal is to identify key features in text to be able to classify key moments in games.

\medskip

As a part of the analysis, Oliveira et al. also discuss the concept of \textit{active time} regarding moments of high intensity. Similar to the work of Pedrassoli covered in Section \ref{subsec:broadcast-medium-tools} this amount of time is relevant to keep the viewer engaged, and often corresponds to narrative moments. The results of this article show great potential with a similar workflow to the one propose in this thesis, although the reliance on multi-modal data causes sensitivity issues, particularly with crowd noise.

\medskip 

Although this paper mentions the possibility of its application to esports, they identify a lack of a curated dataset to perform a pilot study. Addressing this is a key contribution of this thesis, as mentioned for the previous article as well. 

\medskip

Furthermore, this work also aims to identify additional labels, since the focus of Oliveira's work is limited to highlights, rather than the different types of discourse functions which can be found throughout a broadcast.

\subsection{Attempts at Taxonomy and Labeling}

